\subsubsection{Plataforma de processamento}

A plataforma de processamento é responsável pela execução do sistema 
operacional embarcado, pela aquisição dos sinais fisiológicos, pelo 
processamento dos dados coletados, pelo cálculo do escore NEWS2 e pela 
execução dos modelos de aprendizado de máquina empregados na predição de 
deterioração clínica. Dessa forma, sua seleção influencia diretamente o 
desempenho, a confiabilidade e a viabilidade do sistema proposto.

Segundo \citeonline{ozturk2023}, arquiteturas ARM representam a principal 
tendência em sistemas embarcados voltados ao monitoramento em saúde, devido 
à combinação entre baixo consumo energético, ampla disponibilidade de 
ferramentas de desenvolvimento e capacidade de processamento suficiente 
para aplicações de inteligência artificial em borda. Nesse contexto, foram 
selecionadas para avaliação quatro plataformas baseadas em arquitetura ARM, 
todas capazes de executar um sistema operacional Linux embarcado: Raspberry 
Pi 3B+, BeagleBone Black, Orange Pi Zero e STM32MP135F-DK. Todas as placas 
foram adquiridas previamente para a realização dos experimentos previstos 
neste trabalho.

A avaliação das plataformas foi conduzida considerando três critérios 
principais. O primeiro correspondeu ao suporte oficial ao Buildroot, 
ferramenta adotada neste trabalho para a construção de um sistema 
operacional Linux customizado. A disponibilidade de suporte oficial foi 
considerada um fator relevante por reduzir a complexidade de integração, 
aumentar a estabilidade do sistema e facilitar a reprodução do ambiente de 
desenvolvimento.

O segundo critério está relacionado à disponibilidade de interfaces de 
comunicação compatíveis com os sensores empregados no sistema. Em 
particular, verificou-se a presença dos barramentos I\textsuperscript{2}C e 
SPI, amplamente utilizados em aplicações embarcadas para aquisição de 
sinais biomédicos. A disponibilidade dessas interfaces possibilita a 
integração direta com o sensor fotopletismográfico MAX30102 e com o sensor 
de temperatura DS18B20, além de assegurar compatibilidade com futuras 
expansões do hardware.

O terceiro critério considerou a capacidade computacional necessária para 
a execução local dos modelos preditivos. Como o sistema foi concebido para 
operar em computação de borda, avaliou-se a capacidade das plataformas de 
executar inferências em modelos de aprendizado de máquina, como o XGBoost, 
com latência compatível com aplicações de monitoramento contínuo em 
tempo real.

Para subsidiar a comparação entre as plataformas, além das características 
nominais de hardware e do suporte ao Buildroot, foram definidos quatro 
testes práticos a serem executados em cada placa, sob condições 
equivalentes:

\begin{enumerate}
    \item[a)] \textbf{Tempo de boot do sistema operacional}: medição do 
    intervalo entre a energização da placa e a inicialização do serviço 
    principal da aplicação, com sistema gerado pelo Buildroot e 
    configuração mínima de pacotes.
    
    \item \textbf{Latência de inferência de modelo XGBoost de referência}: 
    medição do tempo médio de inferência sobre um vetor de entrada 
    sintético, utilizando um modelo treinado fora da placa e embarcado por 
    meio da biblioteca \textit{treelite}.
    
    \item \textbf{Throughput de leitura via I\textsuperscript{2}C}: medição 
    da taxa efetiva de leitura de amostras do sensor MAX30102 a uma 
    frequência alvo de 100 Hz, verificando a ocorrência de perdas de 
    amostras em janelas de 30 segundos.
    
    \item \textbf{Uso de memória RAM em regime}: medição da memória 
    residente ocupada pela aplicação completa após a estabilização do 
    sistema, incluindo o pipeline de aquisição, o cálculo do NEWS2 e a 
    inferência do modelo preditivo.
\end{enumerate}

As características técnicas e o preço médio de aquisição das plataformas 
avaliadas estão sintetizados na Tabela~\ref{tableComparativoPlataformas}. 
Os preços foram levantados em fornecedores nacionais no período de 
realização deste trabalho e estão sujeitos a variações cambiais e de 
disponibilidade.

\begin{table}[H]
\centering
\caption{Comparação das plataformas de processamento avaliadas}
\label{tableComparativoPlataformas}
\footnotesize
\begin{tabular}{|l|c|c|c|c|c|c|}
\hline
\textbf{Plataforma} &
\textbf{CPU} &
\textbf{RAM} &
\textbf{Buildroot} &
\textbf{I\textsuperscript{2}C/SPI} &
\textbf{Display} &
\textbf{Preço (R\$)} \\
\hline
\textbf{Raspberry Pi 3B+} &
Cortex-A53 4$\times$1,4 GHz &
1 GB &
Oficial &
Sim &
HDMI / DSI &
$\sim$350--500 \\
\hline
BeagleBone Black &
Cortex-A8 1 GHz &
512 MB &
Oficial &
Sim &
HDMI &
$\sim$700--950 \\
\hline
Orange Pi Zero &
Cortex-A7 4$\times$1,2 GHz &
512 MB &
Oficial &
Sim &
HDMI (expansão) &
$\sim$250--400 \\
\hline
STM32MP135F-DK &
Cortex-A7 1 GHz &
512 MB &
Oficial &
Sim &
LCD 4" integrado &
$\sim$800--1200 \\
\hline
\end{tabular}
{\footnotesize\\ \textbf{Fonte: Elaborado pelos autores.}}
\end{table}

Conforme apresentado na Tabela~\ref{tableComparativoPlataformas}, todas as 
plataformas avaliadas disponibilizam interfaces de comunicação compatíveis 
com os sensores empregados neste trabalho, bem como suporte oficial ao 
Buildroot. Entretanto, observou-se que a Raspberry Pi 3B+ apresentou o 
melhor equilíbrio entre os critérios estabelecidos. Além do suporte 
amplamente documentado ao Buildroot e da comunidade ativa que facilita a 
resolução de problemas durante o desenvolvimento, a plataforma oferece o 
maior volume de memória RAM entre as candidatas (1 GB), arquitetura 
quad-core de 64 bits a 1,4 GHz e disponibilidade nativa de interfaces 
I\textsuperscript{2}C, SPI, HDMI e DSI. Tais recursos asseguram capacidade 
computacional suficiente para a execução local de modelos de aprendizado 
de máquina, com folga para o processamento simultâneo da interface 
gráfica.

A BeagleBone Black e a Orange Pi Zero também atenderam aos requisitos de 
comunicação e compatibilidade com o Buildroot, porém apresentaram menor 
capacidade computacional quando comparadas à Raspberry Pi 3B+. A 
BeagleBone Black destaca-se pela presença dos PRUs (\textit{Programmable 
Real-time Units}), que viabilizam aquisição determinística em aplicações 
de tempo real, característica não explorada neste trabalho. A Orange Pi 
Zero apresenta o menor custo entre as candidatas, mas a sua memória RAM 
mais limitada e o ecossistema de documentação menos consolidado tornam o 
processo de desenvolvimento mais oneroso.

A plataforma STM32MP135F-DK, embora baseada em arquitetura ARM Cortex-A7 e 
plenamente capaz de executar um sistema Linux customizado via Buildroot, 
apresenta arquitetura \textit{single-core} e custo de aquisição superior 
ao das demais opções. Tais características reduzem a sua atratividade para 
o cenário específico deste trabalho, no qual a paralelização de tarefas 
entre aquisição contínua, processamento de sinais, inferência do modelo 
preditivo e interface gráfica favorece arquiteturas \textit{multi-core}. 
Adicionalmente, embora o ecossistema OpenSTLinux ofereça integração com 
Buildroot, a sua adoção implicaria uma curva de aprendizado adicional não 
justificada frente aos ganhos esperados.

Dessa forma, a Raspberry Pi 3B+ foi selecionada como plataforma de 
processamento do sistema proposto, constituindo a base de hardware 
utilizada nas etapas subsequentes de desenvolvimento e validação.